\documentclass[11pt]{article}

\usepackage[T1]{fontenc}
\usepackage{lmodern}
\usepackage[margin=1in]{geometry}
\usepackage{amsmath,amssymb,amsthm,mathtools}
\usepackage{microtype}
\usepackage[hidelinks]{hyperref}

\newtheorem{theorem}{Theorem}
\newtheorem{lemma}{Lemma}
\newcommand{\E}{\mathbb E}
\newcommand{\R}{\mathbb R}
\newcommand{\C}{\mathbb C}
\newcommand{\loc}{\mathrm{loc}}
\newcommand{\HS}{\mathrm{HS}}
\DeclareMathOperator{\Unif}{Unif}
\DeclareMathOperator{\Corr}{Corr}
\DeclareMathOperator{\MMSE}{MMSE}
\DeclareMathOperator{\Lift}{Lift}
\DeclareMathOperator{\Id}{Id}
\DeclareMathOperator{\ST}{ST}
\allowdisplaybreaks[1]

\title{Low-Degree Lower Bounds for Overcomplete Tensor Decomposition\thanks{This manuscript was generated with GPT-6 Astra and has not been independently verified by humans.}}
\author{}
\date{}

\begin{document}
\maketitle

\begin{abstract}
We study a form of overcomplete tensor decomposition in the low-degree polynomial framework: given an order-\(k\) tensor with $m$ components $v_1,\dots,v_m$ drawn from $\{\pm 1\}^n$, where $v_1^{\otimes k}$ is scaled up by a constant factor $\tau >1$, the goal is for a polynomial of degree at most \(\ell\) to recover the first entry of $v_1$, or the product of the first two entries if $k$ is even.
For \(k\in\{3,4\}\), degree \(\ell^3 \le n\), and \(m \ge C n^{k/2}\ell^{k/2-1}\) for a large enough constant $C$, we prove that any polynomial of degree at most \(\ell\) has negligible advantage, in mean square error, over the constant-zero estimator.
\end{abstract}

\section{Definitions}

Following Wein~\cite[Sections~1.1--1.2]{Wein2023}, we consider a form of tensor decomposition that fits the low-degree polynomial (LDP) framework: one component has larger norm, and the goal is to estimate its first entry.
Let \([n]=\{1,\ldots,n\}\), and write \(v^A=\prod_{i\in A}v(i)\) for \(A\subseteq[n]\).
Fix \(k\in\{3,4\}\), integers \(n\ge4\) and \(m\ge2\), and \(\tau>1\).
The observation is
\begin{equation}\label{eq:model}
T=\tau v_1^{\otimes k}+\sum_{j=2}^{m}v_j^{\otimes k},
\qquad
v_1,\ldots,v_{m}\stackrel{\mathrm{iid}}{\sim}
\Unif(\{\pm1\}^n).
\end{equation}
Here \(m\) denotes the rank parameter, namely the number of rank-1 terms; the first term has weight \(\tau>1\), and all remaining terms have weight one.
Equivalently, the component vectors are \(\tau^{1/k}v_1,v_2,\ldots,v_m\), so estimating the first entry of the component of larger norm amounts to estimating \(v_1(1)\), up to a known scale.
Define
\begin{equation}\label{eq:observations}
\Omega_k=\{I\subseteq[n]:1\le |I|\le k,\ |I|\equiv k\pmod2\},
\qquad
Y_I=\tau v_1^I+\sum_{j=2}^{m}v_j^I .
\end{equation}
Each tensor entry equals \(Y_I\) for the set \(I\) of coordinates occurring an odd number of times, or, if \(I=\varnothing\), the deterministic constant \(m-1+\tau\).
Conversely, every \(Y_I\), \(I\in\Omega_k\), occurs as a tensor entry.
Thus polynomials in \(T\) and in \(Y=(Y_I)_{I\in\Omega_k}\) generate the same observable polynomial space, with the same degree bound.

The target classes are
\begin{equation}\label{eq:targets}
\mathcal A_3=\{A\subseteq[n]:|A|=1\},
\qquad
\mathcal A_4=\{A\subseteq[n]:|A|\in\{2,4\}\}.
\end{equation}
For \(A\in\mathcal A_k\), put \(x_A=v_1^A\) and \(a=|A|\).
For \(k=3\), taking \(A=\{1\}\) gives the first-entry target \(v_1(1)\).
For \(k=4\), the tensor is unchanged by \(v_1\mapsto-v_1\), so we use sign-invariant targets; the pair target \(v_1(1)v_1(2)\) was suggested by Wein~\cite[Section~2]{Wein2023}.
It is the first entry of the representative \(v_1(2)v_1\), whose second entry is positive.
The theorem also covers all pair and quartet targets in \eqref{eq:targets}.
For a positive integer \(\ell\), let \(\R[T]_{\le\ell}\) be the space of real polynomials of degree at most \(\ell\) in the tensor entries.
Following Wein~\cite[Sections~1.2 and~1.4.1]{Wein2023}, define the \emph{degree-\(\ell\) minimum mean squared error} and the \emph{degree-\(\ell\) maximum correlation} by
\begin{align}
\MMSE_{\le \ell}(x_A;T)
&=\inf_{f\in\R[T]_{\le\ell}}\E[(f(T)-x_A)^2],
\label{eq:mmse}\\
\Corr_{\le \ell}(x_A;T)
&=\sup_{\substack{f\in\R[T]_{\le\ell}\\\E[f(T)^2]>0}}
\frac{\E[x_Af(T)]}{\sqrt{\E[f(T)^2]}} .
\label{eq:corr}
\end{align}
Both optimizations range over all real polynomials in the full tensor entries.
Since \(\E[x_A^2]=1\), the standard identity used in Wein~\cite[Fact~3.1]{Wein2023} gives
\(\MMSE_{\le \ell}(x_A;T)=1-\Corr_{\le \ell}(x_A;T)^2\).

\section{Statement}

\begin{theorem}\label{thm:main}
For every fixed \(\tau>1\), there exists \(C_\tau<\infty\), independent of \(n,m,\ell\) and uniform for \(k\in\{3,4\}\), such that
\begin{equation}\label{eq:hypotheses}
\ell^3\le n,
\qquad
m\ge C_\tau n^{k/2}\ell^{k/2-1}
\end{equation}
imply, for every \(A\in\mathcal A_k\),
\begin{equation}\label{eq:conclusion}
\Corr_{\le \ell}(x_A;T)^2\le n^{-|A|/2},
\qquad
\MMSE_{\le \ell}(x_A;T)\ge1-n^{-|A|/2}.
\end{equation}
\end{theorem}

\section{Proof}

We use Wein's reduction to the Fourier basis of the underlying Rademacher variables~\cite[Sections~1.4.1 and~3.1--3.2]{Wein2023}, expressed below through a moment-matching witness.
Throughout the proof, \(C\) denotes a universal constant that may change between occurrences.
Constants with subscript \(\tau\) depend only on \(\tau\).
Fix an allowed target \(A\), and abbreviate \(\Omega=\Omega_k\).
For uniform \(z\in\{\pm1\}^n\), set
\begin{equation}\label{eq:tilt}
F(t)=\E_z\exp\left(\sum_{I\in\Omega}t_Iz^I\right),
\qquad
\phi_B(t)=
\frac{\E_z[z^B\exp(\sum_{I\in\Omega}t_Iz^I)]}{F(t)}
\quad(\varnothing\ne B\subseteq[n]).
\end{equation}
These functions initially define holomorphic germs at zero, with \(\phi_B(0)=0\).
Let \(K_I(t)\) be the joint cumulant of \((z_i)_{i\in I}\) under this tilt, for \(I\in\Omega\).
For \(k=3\), its coordinates are
\begin{equation}\label{eq:k3}
K_i=\phi_i,\qquad
K_{ijs}
=\phi_{ijs}-\phi_i\phi_{js}-\phi_j\phi_{is}
-\phi_s\phi_{ij}+2\phi_i\phi_j\phi_s .
\end{equation}
For \(k=4\), global spin-flip symmetry makes every odd tilted moment vanish, so
\begin{equation}\label{eq:k4}
K_{ij}=\phi_{ij},\qquad
K_{ijsr}
=\phi_{ijsr}-\phi_{ij}\phi_{sr}
-\phi_{is}\phi_{jr}-\phi_{ir}\phi_{js}.
\end{equation}
All indices within a displayed set are distinct.
Character orthogonality gives \(K(0)=0\) and \(DK(0)=\Id\).
For \(t\in\C^\Omega\), write
\begin{equation}\label{eq:local-norm}
\|t\|_{\loc}=\max_{i\in[n]}\sum_{\substack{I\in\Omega\\i\in I}}|t_I|.
\end{equation}

\begin{lemma}\label{lem:analytic}
There exists \(\rho_\tau>0\), independent of \(n,k,A\), such that
\begin{equation}\label{eq:response}
\mathcal C_A(q)=\phi_A(\tau K^{-1}(q))
\end{equation}
is holomorphic on a neighborhood of the closed ball
\(\|q\|_{\loc}\le\rho_\tau\) and satisfies \(|\mathcal C_A(q)|\le1\) there.
\end{lemma}

\begin{proof}
Write \(K(t)=\sum_{r\ge1}T_r(t,\ldots,t)\).
Differentiating tilted cumulants gives the symmetric multilinear coefficients
\begin{equation}\label{eq:multilinear}
T_{r,I}(w_1,\ldots,w_r)
=\frac1{r!}\sum_{J_1,\ldots,J_r\in\Omega}
\kappa_0\big((z_i)_{i\in I},z^{J_1},\ldots,z^{J_r}\big)
\prod_{b=1}^r(w_b)_{J_b}.
\end{equation}
We use the dependency-graph inequality~\cite[Eq.~(47)]{FMN2015}
\begin{equation}\label{eq:tree-bound}
|\kappa(X_1,\ldots,X_s)|\le2^{s-1}\ST(G)
\end{equation}
for variables bounded in absolute value by one with dependency graph \(G\), where \(\ST(G)\) is the number of spanning trees.
For \eqref{eq:multilinear}, join character occurrences whose supports intersect.
This is a dependency graph, including when occurrences repeat.

Fix an output coordinate \(i\) and output size \(p\in\{k-2,k\}\).
Root a spanning tree at the marked spin \(z_i\).
An interaction child \(b\) attached to a parent support \(B\) contributes at most
\[
\sum_{J:J\cap B\ne\varnothing}|(w_b)_J|
\le |B|\|w_b\|_{\loc}\le k\|w_b\|_{\loc}.
\]
Each of the other \(p-1\) marked spin children has at most \(k\) coordinate choices.
These subtree bounds are uniform in the parent's support.
Summing ordered choices of marked coordinates overcounts the output sets.
Thus Cayley's formula and \eqref{eq:tree-bound} give
\begin{equation}\label{eq:anchored-tree}
\sum_{\substack{I\ni i\\|I|=p}}
|T_{r,I}(w_1,\ldots,w_r)|
\le
\frac{(2k)^{r+p-1}(r+p)^{r+p-2}}{r!}
\prod_{b=1}^r\|w_b\|_{\loc}.
\end{equation}
Since \(k,p\le4\), the sum of the two numerical prefactors is at most \(C_0^r\) for a universal \(C_0\).
Consequently
\begin{equation}\label{eq:mixed-bound}
\|T_r(w_1,\ldots,w_r)\|_{\loc}
\le C_0^r\prod_{b=1}^r\|w_b\|_{\loc}.
\end{equation}

The resulting series converges on a dimension-independent complex ball.
To identify it with the actual moment ratios, let \(L_r\) be the homogeneous degree-\(r\) part of the formal logarithm of \(F\).
Here \(L_1=0\).
Rooting the interaction-only tree at its first interaction and using
\(\sum_I|t_I|\le n\|t\|_{\loc}\) gives, for \(r\ge2\),
\begin{equation}\label{eq:log-bound}
|L_r(t)|
\le
\frac{2^{r-1}r^{r-2}}{r!}
n\|t\|_{\loc}(k\|t\|_{\loc})^{r-1}
\le nC^r\|t\|_{\loc}^r.
\end{equation}
The logarithmic series converges on a universal ball.
The germ identity \(F=\exp(\sum_rL_r)\) extends throughout that connected ball, so \(F\) has no zeros there.
The series for \(K\) therefore agrees with \eqref{eq:k3}--\eqref{eq:k4}.

Set \(H=K-\Id\).
Since \(T_1=\Id\), the mixed bound \eqref{eq:mixed-bound} implies
\begin{equation}\label{eq:derivative-bound}
\|DH(t)\|_{\loc\to\loc}
\le\sum_{r\ge2}rC_0^r\|t\|_{\loc}^{r-1}.
\end{equation}
Choose a universal \(\eta>0\) so that the right-hand side is at most \(1/2\) on the closed \(\eta\)-ball, inside the analytic domain.
For \(\|q\|_{\loc}<\eta/2\), the map \(t\mapsto q-H(t)\) contracts that ball.
Its holomorphic iterates converge uniformly to the inverse, and
\begin{equation}\label{eq:inverse-bound}
\|K^{-1}(q)\|_{\loc}\le2\|q\|_{\loc},
\qquad
\|K(t)\|_{\loc}\le\tfrac32\|t\|_{\loc}
\quad(\|t\|_{\loc}\le\eta).
\end{equation}
For singleton targets when \(k=3\) and pair targets when \(k=4\), \(\phi_A\) is a coordinate of \(K\).
For a quartet,
\[
\phi_{ijsr}
=K_{ijsr}+K_{ij}K_{sr}
+K_{is}K_{jr}+K_{ir}K_{js}.
\]
Shrinking \(\eta\), and then taking a sufficiently small \(\rho_\tau<\eta/(4\tau)\), proves the asserted holomorphy and bound for \(\mathcal C_A\).
\end{proof}

For a multiset \(S\) of elements of \(\Omega\), write
\[
|S|=\sum_I S_I,\qquad S!=\prod_I S_I!,\qquad
q^S=\prod_Iq_I^{S_I},
\]
and let \(\triangle S\) be the symmetric difference of its elements, with multiplicity.
Expand
\begin{equation}\label{eq:energy-definition}
\mathcal C_A(q)=\sum_{d\ge1}\mathcal C_{A,d}(q),
\qquad
\mathcal C_{A,d}(q)=\sum_{|S|=d}\frac{c_{A,S}}{S!}q^S,
\qquad
\mathcal E_{A,d}=\sum_{|S|=d}\frac{c_{A,S}^2}{S!}.
\end{equation}
The coefficients are real.

\begin{lemma}\label{lem:coefficients}
For all \(d\ge1\),
\begin{equation}\label{eq:energy-bound}
\mathcal E_{A,d}
\le n^{-a/2}
\bigl(B_\tau n^{k/2}d^{k/2-1}\bigr)^d
\exp(4d^4/n),
\end{equation}
where \(B_\tau\) is independent of \(n,k,A,d\).
\end{lemma}

\begin{proof}
Coordinate-sign changes commute with \(K\) and its unique local inverse, and multiply the response by the sign character indexed by \(A\).
Hence
\begin{equation}\label{eq:parity}
c_{A,S}=0\quad\text{unless}\quad\triangle S=A.
\end{equation}
Equivalently, \(S\) together with \(A\) forms an \emph{even cover}: every coordinate occurs an even number of times, counting multiplicities.
This is the even-cover restriction used by Wein~\cite[Lemma~3.16]{Wein2023}, with the singleton target there replaced by \(A\).
In particular, contributing degrees satisfy \(kd\equiv a\pmod2\).
For \(k=3\) this requires odd \(d\); for \(k=4\) it imposes no degree restriction.

For a fixed multiset \(S\), set \(s_i=\sum_{I\ni i}S_I\) and \(M_I=\max_{i\in I}s_i\).
On the active variables, choose radii \(r_I=\rho_\tau S_I/M_I\).
Since
\(\sum_{I\ni i}r_I\le(\rho_\tau/s_i)\sum_{I\ni i}S_I=\rho_\tau\)
for every occurring coordinate \(i\), this polydisc lies in the ball of Lemma~\ref{lem:analytic}.
Cauchy's estimate yields
\begin{equation}\label{eq:cauchy}
|c_{A,S}|
\le S!\prod_I r_I^{-S_I}
\le\rho_\tau^{-d}\prod_I M_I^{S_I}
\le\rho_\tau^{-d}\prod_{i:s_i>0}s_i^{s_i}.
\end{equation}
Here we used \(S_I!\le S_I^{S_I}\) and \(M_I\le\prod_{i\in I}s_i\).

We refine Wein's counting of even covers~\cite[Lemma~3.17]{Wein2023} by keeping track of coordinate multiplicities.
Passing from multisets to ordered tuples gives
\begin{equation}\label{eq:ordered-count}
\mathcal E_{A,d}
\le\frac{\rho_\tau^{-2d}}{d!}
\sum_{\substack{I_1,\ldots,I_d\in\Omega\\\triangle_b I_b=A}}
\prod_{i:s_i>0}s_i^{2s_i}.
\end{equation}
Fix one of the at most \(2^d\) choices of input sizes in \(\{k-2,k\}\).
Let \(s\le kd\) be its number of incidence slots.
A contributing pattern has \(s\ge a\) and \(s\equiv a\pmod2\); put \(L=(s-a)/2\).
The multiplicities can be written as
\[
s_i=1+2b_i\quad(i\in A),\qquad
s_j=2+2b_j\quad(j\notin A\text{ occurring}).
\]
If \(h\) nonroot coordinates occur and \(e=\sum_i b_i\), then \(h=L-e\).
There are at most \(n^h/h!\) supports and
\(\binom{e+h+a-1}{e}=\binom{L+a-1}{e}\)
weak compositions of \(e\).

For fixed multiplicities, there are \(s!/\prod_i s_i!\) assignments to the labeled slots.
Allowing repetitions within an input set only enlarges the count.
If \(u=b+2b'\le4d\), where \(b\in\{1,2\}\), then
\begin{equation}\label{eq:multiplicity-bound}
\frac{u^{2u}}{u!}\le \exp(u)u^u
\le(16\exp(1))^u d^{2b'}.
\end{equation}
Indeed \(u^u\le(4d)^{2b'}u^2\), \(4^{2b'}\le4^u\), and \(u^2\le4^u\).
Thus the fixed size pattern contributes at most
\begin{equation}\label{eq:collision-sum}
\frac{C_\tau^d s!}{d!}
\sum_{e=0}^{L}\frac{n^{L-e}}{(L-e)!}
\binom{L+a-1}{e}d^{2e}.
\end{equation}
If the support size exceeds the available coordinates, its actual contribution is zero, so the relaxed support bound remains valid.

After factoring out \(n^L/L!\), use
\begin{equation}\label{eq:collision-exponential}
\frac{L!}{(L-e)!}\binom{L+a-1}{e}d^{2e}
\le\frac{[L(L+a-1)d^2]^e}{e!}
\le\frac{(4d^4)^e}{e!}.
\end{equation}
The last inequality follows from \(L\le(3d-1)/2\), \(2d-1\), or \(2d-2\) for \((k,a)=(3,1),(4,2),(4,4)\), respectively.
When \(L=0\), only \(e=0\) occurs and the expression is one.
Summing in \eqref{eq:collision-sum} therefore gives
\begin{equation}\label{eq:fixed-size-bound}
\frac{C_\tau^d s!n^L}{d!L!}\exp(4d^4/n).
\end{equation}

Increasing \(s\) by two multiplies \(s!n^{(s-a)/2}/((s-a)/2)!\) by
\(n(s+2)(s+1)/(L+1)\ge1\).
Since \(kd\equiv a\pmod2\), the maximum occurs at \(s=kd\).
Put \(L_*=(kd-a)/2\).
Except for \((k,a,d)=(4,4,1)\), which is immediate, \(L_*\ge d\), and elementary factorial bounds give
\[
\frac{(kd)!}{d!L_*!}
\le C^d d^{kd-d-L_*}
=C^d d^{(k/2-1)d+a/2}
\le C^d d^{(k/2-1)d}.
\]
In the last inequality the constant is enlarged using \(d^{a/2}\le d^2\le4^d\).
Now \(n^{L_*}=n^{-a/2}n^{kd/2}\); summing the at most \(2^d\) size patterns proves \eqref{eq:energy-bound}.
\end{proof}

Define the auxiliary index sets
\[
\Lambda_3=\{B\subseteq[n]:1\le|B|\le3\},
\qquad
\Lambda_4=\Omega_4.
\]
Following Wein's use of Fourier monomials with distinct component indices~\cite[Section~3.4]{Wein2023}, we average over distinct background components.
For a multiset \(R=(B_1,\ldots,B_r)\) from \(\Lambda_k\), with \(r\le m-1\), let
\begin{equation}\label{eq:fourier-witnesses}
U_R=\frac1{(m-1)^{\underline r}}
\sum_{\substack{j_1,\ldots,j_r\in\{2,\ldots,m\}\\\text{all distinct}}}
\prod_{b=1}^r v_{j_b}^{B_b},
\qquad
(m-1)^{\underline r}=(m-1)(m-2)\cdots(m-r).
\end{equation}
Set \(U_\varnothing=1\).
For a polynomial \(P(u)=\sum_Rp_Ru^R\) of degree at most \(m-1\), define
\(\Lift(P)=\sum_Rp_RU_R\).
The auxiliary pair characters used for \(k=3\) are functions of the underlying Rademacher variables, not additional observations.

Let \(b_k(u)\) be obtained from \eqref{eq:k3} or \eqref{eq:k4} by replacing each \(\phi_B\) with the independent letter \(u_B\).
Then
\begin{equation}\label{eq:substitution}
b_k(\phi(t))=K(t),\qquad
\deg b_k\le h_k,\qquad h_3=3,\quad h_4=2.
\end{equation}
When \(m-1\ge h_k\ell\), set
\begin{equation}\label{eq:witness-definition}
P_{A,\ell}(u)=\sum_{d=1}^\ell\mathcal C_{A,d}(b_k(u)),
\qquad
W_{A,\ell}=\Lift(P_{A,\ell}).
\end{equation}

\begin{lemma}\label{lem:lift}
Suppose \(m-1\ge6\ell\).
Then \(W_{A,\ell}\) satisfies
\begin{equation}\label{eq:moment-matching}
\E[W_{A,\ell}f(Y)]=\E[x_Af(Y)]
\qquad(\deg f\le \ell),
\end{equation}
and, for \(1\le d\le \ell\),
\begin{equation}\label{eq:lift-bound}
\bigl\|\Lift[\mathcal C_{A,d}(b_k(u))]\bigr\|_2
\le6^d\sqrt{\frac{\mathcal E_{A,d}}{(m-1)^{\underline d}}}.
\end{equation}
\end{lemma}

\begin{proof}
Products of the underlying independent Rademacher variables form the standard Fourier basis for functions on the hypercube.
Every distinct basis monomial in \(U_R\) occurs \(R!\) times, where \(R!=\prod_B R_B!\), and there are \((m-1)^{\underline r}/R!\) such monomials.
Different multisets have disjoint Fourier supports.
Consequently
\begin{equation}\label{eq:orthogonality}
\langle U_R,U_{R'}\rangle=0\quad(R\ne R'),
\qquad
\|U_R\|_2^2=\frac{R!}{(m-1)^{\underline{|R|}}}.
\end{equation}
Independence gives
\begin{align}
\E[U_Re^{\langle t,Y\rangle}]
&=F(\tau t)F(t)^{m-1}\phi(t)^R,
\label{eq:tilted-witness}\\
\E[x_Ae^{\langle t,Y\rangle}]
&=F(\tau t)F(t)^{m-1}\phi_A(\tau t).
\label{eq:tilted-target}
\end{align}
Since \(K(0)=0\) and \(\mathcal C_A(K(t))=\phi_A(\tau t)\), the polynomial in \eqref{eq:witness-definition} satisfies
\[
P_{A,\ell}(\phi(t))=\phi_A(\tau t)+O_{\deg}(\ell+1).
\]
Its degree in the auxiliary variables is at most \(h_k\ell\le3\ell\le m-1\).
Multiplying the formal identity by \(F(\tau t)F(t)^{m-1}\) and comparing Taylor coefficients in \eqref{eq:tilted-witness}--\eqref{eq:tilted-target} proves \eqref{eq:moment-matching}.

For the norm bound, let \(V=\R^\Omega\) and \(H=\R^{\Lambda_k}\), with standard orthonormal bases.
Represent the homogeneous parts of \(b_k\) by symmetric tensor maps \(Q_j:V\to H^{\otimes j}\):
\[
(b_k)_I^{[j]}(u)=\langle Q_je_I,u^{\otimes j}\rangle.
\]
The linear map is the basis inclusion.
For \(k=3\), a triple output has six ordered singleton--pair coefficients \(-1/2\) in \(Q_2\), and six ordered three-singleton coefficients \(1/3\) in \(Q_3\).
For \(k=4\), a quartet output has six ordered complementary-pair coefficients \(-1/2\) in \(Q_2\).
Within each map, distinct output sets have disjoint tensor supports, because the union of an input tuple identifies its output.
Thus
\begin{equation}\label{eq:map-norms}
\begin{array}{c|ccc}
&\|Q_1\|&\|Q_2\|&\|Q_3\|\\ \hline
k=3&1&\sqrt{3/2}&\sqrt{2/3}\\
k=4&1&\sqrt{3/2}&0
\end{array}
\qquad\text{and}\qquad
\sum_{j=1}^{h_k}\|Q_j\|\le6.
\end{equation}

Write \(\mathcal C_{A,d}(q)=\langle M_d,q^{\otimes d}\rangle\), with \(M_d\) symmetric.
Its entry on an ordered tuple of multiset \(S\) is \(c_{A,S}/d!\), so
\begin{equation}\label{eq:tensor-normalization}
d!\|M_d\|_{\HS}^2=\mathcal E_{A,d}.
\end{equation}
For any symmetric order-\(s\) tensor \(M\), \eqref{eq:orthogonality} implies the exact identity
\begin{equation}\label{eq:lift-normalization}
\bigl\|\Lift\langle M,u^{\otimes s}\rangle\bigr\|_2^2
=\frac{s!}{(m-1)^{\underline s}}\|M\|_{\HS}^2.
\end{equation}
For one substitution choice \((j_1,\ldots,j_d)\), the resulting tensor has order \(s=\sum_bj_b\in[d,h_kd]\) and norm at most
\(\prod_b\|Q_{j_b}\|\cdot\|M_d\|_{\HS}\).
Orthogonal symmetrization leaves its polynomial unchanged and cannot increase the norm.
When \(m-1\ge6\ell\), the successive ratio
\[
\frac{(s+1)!/(m-1)^{\underline{s+1}}}{s!/(m-1)^{\underline s}}
=\frac{s+1}{m-1-s}\le1
\qquad(d\le s<3d)
\]
shows that \(s!/(m-1)^{\underline s}\le d!/(m-1)^{\underline d}\) throughout the needed range.
Using \eqref{eq:tensor-normalization}--\eqref{eq:lift-normalization} and summing all substitution choices by the triangle inequality gives
\[
\bigl\|\Lift[\mathcal C_{A,d}(b_k(u))]\bigr\|_2
\le
\left(\sum_{j=1}^{h_k}\|Q_j\|\right)^d
\sqrt{\frac{\mathcal E_{A,d}}{(m-1)^{\underline d}}},
\]
which proves \eqref{eq:lift-bound}.
\end{proof}

\begin{proof}[Proof of Theorem~\ref{thm:main}]
For \(d\le \ell\), the assumption \(\ell^3\le n\) gives \(d^4/n\le d\).
Absorbing the exponential in Lemma~\ref{lem:coefficients} into its constant, we obtain
\begin{equation}\label{eq:absorbed-bound}
\mathcal E_{A,d}
\le n^{-a/2}\bigl(B_\tau n^{k/2}d^{k/2-1}\bigr)^d.
\end{equation}
Choose \(C_\tau\ge\max\{8,288B_\tau\}\).
Under \eqref{eq:hypotheses}, \(n^{k/2}\ell^{k/2-1}\ge \ell\), so \(m-1\ge6\ell\) and \(m-\ell\ge m/2\).
By Lemma~\ref{lem:lift}, the triangle inequality across degrees, and
\((m-1)^{\underline d}\ge(m-\ell)^d\),
\begin{equation}\label{eq:geometric-bound}
\|W_{A,\ell}\|_2
\le n^{-a/4}
\sum_{d=1}^\ell
\left(
\frac{36B_\tau n^{k/2}d^{k/2-1}}{m-\ell}
\right)^{d/2}.
\end{equation}
Each base is at most
\[
\frac{72B_\tau n^{k/2}\ell^{k/2-1}}{m}\le\frac14,
\]
and hence \(\|W_{A,\ell}\|_2\le n^{-a/4}\sum_{d\ge1}2^{-d}=n^{-a/4}\).
For every polynomial \(f\) of degree at most \(\ell\), moment matching yields
\[
\E[(f(Y)-x_A)^2]
=1-\|W_{A,\ell}\|_2^2+\E[(f(Y)-W_{A,\ell})^2]
\ge1-n^{-a/2}.
\]
Alternatively, Cauchy--Schwarz applied to
\(\E[x_Af(Y)]=\E[W_{A,\ell}f(Y)]\)
gives the squared-correlation bound directly.
The equivalence of the \(T\)- and \(Y\)-polynomial spaces completes the proof.
\end{proof}

\begin{thebibliography}{9}

\bibitem{Wein2023}
Alexander S.~Wein.
\newblock Average-case complexity of tensor decomposition for low-degree polynomials.
\newblock In \emph{Proceedings of the 55th Annual ACM Symposium on Theory of Computing (STOC 2023)}, pp.~1685--1698, 2023.
\newblock \href{https://doi.org/10.1145/3564246.3585232}{doi:10.1145/3564246.3585232}.
\newblock \href{https://arxiv.org/abs/2211.05274v2}{arXiv:2211.05274v2}.

\bibitem{FMN2015}
Valentin F\'eray, Pierre-Lo\"ic M\'eliot, and Ashkan Nikeghbali.
\newblock \emph{Mod-\(\phi\) convergence, I: Normality zones and precise deviations}.
\newblock \href{https://arxiv.org/abs/1304.2934v4}{arXiv:1304.2934v4}, 2015.

\end{thebibliography}

\end{document}
